\documentclass{article}
\usepackage[utf8]{inputenc}
\usepackage{amsmath,amssymb}
\usepackage[most]{tcolorbox}
\usepackage{geometry}
\geometry{margin=2.5cm}

\newcommand{\step}[1]{\par\noindent\hangindent=3.5em\hangafter=1%
  \makebox[3.5em][l]{\textbf{#1.}}}
\newcommand{\steptag}[1]{\hfill{\small\textsf{[#1]}}}

\newtcolorbox{proofbox}[1]{
  colback=gray!5, colframe=gray!60,
  fonttitle=\bfseries, title={#1},
  breakable, enhanced,
  left=4pt, right=4pt, top=4pt, bottom=4pt,
  before skip=10pt, after skip=10pt
}

\begin{document}

\begin{proofbox}{The near-fixed algebra A=Im P, with unit 1 and inherited ...}
\step{1} The near-fixed algebra A=Im P, with unit 1 and inherited cone A\_+ = A cap B(H)\_+, is an order-unit space, and its order-unit norm ||a||\_ou = inf\{t>0: -t1<=a<=t1\} coincides exactly with the operator norm on A. \steptag{validated} \\[6pt]
\step{1.1} 1 in A and A is a real linear subspace of V=B(H)\_sa. By GT-Pprops, P(1)=1, so 1=P(1) in Im P=A (def-near-fixed-algebra), giving 1 in A. By GT-Pprops, P is real-linear on B(H)\_sa, so its range A=Im P is a real linear subspace of V=B(H)\_sa (image of a real-linear map is a real subspace). Thus A is a real linear subspace of V containing the unit 1. Uses GT-Pprops, def-near-fixed-algebra. \steptag{validated} \\[4pt]
\step{1.2} 1 is an order unit for (A,A\_+): (i) cone-membership 1 in A\_+ holds by node 1.2.1; (ii) dominance — for every a in A, -||a||1 <= a <= ||a||1 in V, since V=B(H)\_sa is a unital JB algebra (GT-bhsa-jc + GT-jc-is-jb) whose operator norm equals its order-unit norm (GT-jb-orderunit-3310), and the order on A is inherited from V (def-near-fixed-algebra, node 1.1). Both clauses give: 1 is an order unit for A. \steptag{validated} \\[6pt]
\step{1.2.1} 1 in A\_+: 1 in A since P(1)=1 (GT-Pprops) and A=Im P (def-near-fixed-algebra). 1 in V\_+ since V=B(H)\_sa is a complete order-unit space with order unit 1 (GT-jb-orderunit-3310) and an order unit lies in the cone (GT-orderunit-def: e in A\textasciicircum{}+). Hence 1 in A cap V\_+ = A\_+ (def-near-fixed-algebra, cone inherited exactly). \steptag{validated} \\[4pt]
\step{1.3} A\_+=A cap V\_+ is Archimedean. We show (A,A\_+,1) satisfies the Archimedean axiom of GT-orderunit-def: if na<=\_A 1 for all n in N then a<=\_A 0. Since the order on A is inherited (def-near-fixed-algebra; A subseteq V by node 1.1), na<=\_A 1 is equivalent to na<=\_V 1 in V. By GT-jb-orderunit-3310 V is a complete order-unit space, hence (GT-orderunit-def) V is itself Archimedean: na<=\_V 1 for all n in N implies a<=\_V 0, i.e. a in V\_+ with -a in V\_+... more precisely a<=\_V 0 means -a in V\_+. Since a in A, a<=\_V 0 gives -a in A cap V\_+=A\_+, i.e. a<=\_A 0. (Equivalently, the same conclusion follows because V\_+ is norm-closed: GT-jb-cone-closed-335 gives (tilde-V)\_+ closed in tilde-V, and V\_+=V cap (tilde-V)\_+ is closed in V; if a+eps 1 in V\_+ for all eps>0 then letting eps->0, a=lim(a+eps 1) in V\_+. Both routes use only inherited order.) Hence A is Archimedean. Uses node 1.1, node 1.2, GT-jb-orderunit-3310, GT-orderunit-def, GT-jb-cone-closed-335, def-near-fixed-algebra. \steptag{validated} \\[4pt]
\step{1.4} (A,A\_+,1) is an order-unit space. By GT-orderunit-def, a real vector space with a proper convex cone and a distinguished element is an order-unit space iff that element is an order unit AND the space is Archimedean. We verify the hypotheses: (i) A is a real vector space with the inherited cone A\_+=A cap V\_+, which is a proper convex cone: convex and closed under R\_+ scaling since V\_+ is (intersection of the subspace A with the convex cone V\_+), and proper (A\_+ cap (-A\_+) = A cap V\_+ cap (-V\_+) = A cap \{0\} = \{0\}) because V\_+ is proper in V (V an order-unit space, GT-jb-orderunit-3310 + GT-orderunit-def); (ii) 1 in A\_+ and 1 is an order unit for A by node 1.2; (iii) A is Archimedean by node 1.3. By GT-orderunit-def these are exactly the defining conditions, so (A,A\_+,1) is an order-unit space (def-order-unit-space). Uses node 1.1, node 1.2, node 1.3, GT-orderunit-def, GT-jb-orderunit-3310, def-order-unit-space, def-near-fixed-algebra. \steptag{validated} \\[4pt]
\step{1.5} ||a||\_\{ou,A\} = ||a||\_\{B(H)\} for all a in A. By node 1.4, (A,A\_+,1) is an order-unit space, so by GT-orderunit-def its order-unit norm is ||a||\_\{ou,A\}=inf\{t>0: -t1<=\_A a<=\_A t1\}. By node 1.1 A subseteq V and the order is inherited (def-near-fixed-algebra), so for a in A and t>0, -t1<=\_A a<=\_A t1 iff -t1<=\_V a<=\_V t1; hence the two infima are over the SAME set of t, giving ||a||\_\{ou,A\}=inf\{t>0:-t1<=\_V a<=\_V t1\}=||a||\_\{ou,V\}. By GT-jb-orderunit-3310, V=B(H)\_sa is a unital JB algebra whose order norm equals the given operator norm, i.e. ||b||\_\{ou,V\}=||b||\_\{B(H)\} for all b in V. Therefore ||a||\_\{ou,A\}=||a||\_\{ou,V\}=||a||\_\{B(H)\} for all a in A. Uses node 1.1, node 1.4, GT-jb-orderunit-3310, GT-orderunit-def, def-order-unit-space, def-near-fixed-algebra. \steptag{validated} \\[4pt]
\end{proofbox}

\end{document}
